\documentclass[11pt]{article}
\usepackage[T1]{fontenc}
\usepackage{amsmath,amssymb,amsthm,mathtools}
\usepackage[margin=1in]{geometry}
\usepackage[colorlinks=true,linkcolor=blue,urlcolor=blue]{hyperref}
\newtheorem{proposition}{Proposition}[section]
\newtheorem{remark}[proposition]{Remark}
\DeclareMathOperator{\diag}{diag}
\DeclareMathOperator{\HeunG}{HeunG}
\newcommand{\Id}{\mathbf 1}
\newcommand{\cI}{\mathcal I}
\newcommand{\Poly}{\mathcal P}
\title{Differential Realization of the Krawtchouk Algebra\\and Its Heun Reductions}
\author{Reference note}
\date{18 October 2026}
\begin{document}
\maketitle
\begin{abstract}
A separate record of the finite polynomial realization of the Krawtchouk
operators, the cubic phase-modified transform, and the resulting Heun
polynomial problems. All conventions use
$\sigma=\diag((-1)^j)$ and $\mathcal I=\diag(i^j)$.
The note records exact finite-dimensional identities, not a proof of a
large-$N$ asymptotic. The current asymptotic strategy through a positive
Jacobi eigenvector and binomially weighted norms does not require the
Heun differential equations; they are retained here for reference.
\end{abstract}
\tableofcontents

\section{Conventions and the finite polynomial model}
Fix an integer $N\geq 1$ and use indices $0,\ldots,N$. Identify
$e_j$ with $z^j$ in $\Poly_N=\{f:\deg f\leq N\}$.
Write $D=d/dz$, $\mathcal E=zD$, and use $\Id$ for the identity.
The coefficient inner product is
\[
 \left\langle\sum f_jz^j,\sum g_jz^j\right\rangle_{\rm coeff}
 =\sum_{j=0}^N\overline{f_j}g_j.
\]
Adjoint and symmetry statements about matrices refer to this inner
product, not to an unspecified integral inner product on a real interval.

Our phase and parity conventions are
\begin{equation}\label{eq:phases}
 \sigma=\diag((-1)^j),\qquad \cI=\diag(i^j),\qquad i^2=-1.
\end{equation}
On polynomials,
\[
 (\sigma f)(z)=f(-z),\qquad (\cI f)(z)=f(iz),\qquad
 \cI^2=\sigma,\quad \cI^{-1}=\overline{\cI}.
\]
Dependence of these matrices on $N$ is suppressed.

\section{Raising, lowering, and diagonal generators}
The finite generators are
\[
 R_Ne_j=(N-j)e_{j+1},\quad L_Ne_j=je_{j-1},\quad
 \Lambda_Ne_j=(N-2j)e_j.
\]
Endpoint terms are omitted. Their differential realizations are
\begin{equation}\label{eq:generators}
 L_N=D,\qquad R_N=Nz-z^2D,\qquad \Lambda_N=N-2zD.
\end{equation}
Consequently,
\[
 A_N=R_N+L_N=(1-z^2)D+Nz,
\]
and
\[
 [L_N,R_N]=\Lambda_N,\quad
 [\Lambda_N,R_N]=-2R_N,\quad
 [\Lambda_N,L_N]=2L_N.
\]
All these operators preserve $\Poly_N$. Preservation depends on the
cancellation of the apparent degree-$N+1$ term in $R_Nz^N$.

\section{The Krawtchouk transform and the symmetric matrix}
Define
\[
 (\Phi_N)_{kj}=[z^k](1+z)^{N-j}(1-z)^j,\qquad
 B_N=\diag\binom Nj,\qquad S_N=\Phi_NB_N.
\]
The transform has the weighted-composition realization
\begin{equation}\label{eq:Phi}
 (\Phi_N f)(z)=(1+z)^N
 f\!\left(\frac{1-z}{1+z}\right).
\end{equation}
This is a polynomial identity, with apparent poles canceled by the
prefactor. It gives $\Phi_N^2=2^N\Id$.
Weighted Krawtchouk duality implies $S_N^T=S_N$.

For clarity, $S_N$ is not being claimed to be a fixed low-order
local differential operator. Define on $\Poly_N$
\[
 \mathscr B_N\!\left(\sum f_jz^j\right)
   =\sum\binom Nj f_jz^j.
\]
Then its exact realization is
\begin{equation}\label{eq:S}
 (S_Nf)(z)=(1+z)^N
 (\mathscr B_N f)\!\left(\frac{1-z}{1+z}\right).
\end{equation}
If a differential expression is desired, $\mathscr B_N=b_N(\mathcal E)$,
where $b_N$ is the interpolation polynomial of degree at most $N$ with
$b_N(j)=\binom Nj$ for $0\leq j\leq N$. This is a finite-space identity.
A useful exact relation, independent of differential realization, is
\begin{equation}\label{eq:normbridge}
 S_NB_N^{-1}S_N=2^NB_N.
\end{equation}

\section{The cubic phase-modified transform}
Put $C_N=\cI\Phi_N$. Then
\begin{equation}\label{eq:C}
 (C_Nf)(z)=(1+iz)^N
 f\!\left(\frac{1-iz}{1+iz}\right).
\end{equation}
Iteration gives
\begin{align}
 (C_N^2f)(z)&=(1+i)^N(1+z)^N
 f\!\left(-i\frac{1-z}{1+z}\right),\label{eq:C2}\\
 C_N^3&=[2(1+i)]^N\Id.\label{eq:cubic}
\end{align}
For example, the degree-one weighted-composition matrix is
$\begin{pmatrix}-i&1\\i&1\end{pmatrix}$, whose cube is
$2(1+i)\Id$; its degree-$N$ homogeneous action gives \eqref{eq:cubic}.
Useful conjugation formulas are
\begin{align}
 C_N^2R_NC_N^{-2}&=\frac i2(R_N-L_N-\Lambda_N),\\
 C_N^2L_NC_N^{-2}&=\frac i2(R_N-L_N+\Lambda_N),\\
 C_N^2\Lambda_NC_N^{-2}&=R_N+L_N.
\end{align}
They provide an algebraic route to the differential operator below.

\section{The skew operator and its positive Jacobi form}
The finite skew operator is
\begin{align}
 T_N&=(\Lambda_N-N\Id)R_N+(\Lambda_N+N\Id)L_N,\\
 T_Ne_j&=2j(N-j+1)e_{j-1}-2(j+1)(N-j)e_{j+1}.
\end{align}
Its differential realization is
\begin{equation}\label{eq:T}
 \boxed{T_N=2\left[z(z^2-1)D^2+
 \bigl(N+(2-N)z^2\bigr)D-Nz\right].}
\end{equation}
It preserves $\Poly_N$ and satisfies
\[
 T_N^T=-T_N,\qquad T_NS_N=-S_NT_N.
\]
Thus $[T_N,S_N^2]=0$.

The phase conjugate
\begin{equation}\label{eq:J}
 \mathsf J_N=-i\cI^{-1}T_N\cI
\end{equation}
is real symmetric, with off-diagonal entries
$2(j+1)(N-j)>0$. Its polynomial expression is
\begin{equation}
 \mathsf J_N=2\left[-z(1+z^2)D^2+
 \bigl(N+(N-2)z^2\bigr)D+Nz\right].
\end{equation}
Let $\tau_N$ denote its largest eigenvalue and $v_N$ its positive
Perron eigenvector. A convenient normalization for this note is
$(v_N)_0=1$. Set $f_N(z)=\sum i^j(v_N)_jz^j$.
Then
\begin{equation}\label{eq:f-eig}
 T_Nf_N=i\tau_Nf_N,\qquad f_N(0)=1.
\end{equation}
Explicitly,
\[
 2z(z^2-1)f_N''+2\bigl(N+(2-N)z^2\bigr)f_N'
 -(2Nz+i\tau_N)f_N=0.
\]

\section{The cubic-conjugated, parity-preserving operator}
Define
\[
 F_N=C_N^2f_N,\qquad
 \mathscr H_N=\frac1iC_N^2T_NC_N^{-2}.
\]
(The symbol $\mathscr H_N$ here denotes this finite transformed operator,
not the continuum harmonic oscillator.) Then
\begin{equation}\label{eq:H}
 \boxed{\begin{aligned}
 \mathscr H_N={}&(z^4-1)D^2
 -2\bigl((N-1)z^3+(N+1)z\bigr)D\\
 &+N(N-1)z^2+N(N+1),
 \end{aligned}}
\end{equation}
and $\mathscr H_NF_N=\tau_NF_N$.
This operator preserves even and odd polynomials separately.

For the Perron mode, $F_N$ is even. To see the selection algebraically,
let $\mathcal Rf(z)=z^Nf(1/z)$ be coefficient reversal.
The Jacobi matrix commutes with $\mathcal R$, so positivity and
simplicity give $\mathcal Rv_N=v_N$. Directly,
\[
 C_N^{-2}\sigma C_N^2=\cI\mathcal R\cI^{-1}.
\]
It follows that $\sigma F_N=F_N$.

Write $N=2M+d$ with $d\in\{0,1\}$ and put
\begin{equation}\label{eq:g}
 F_N(z)=g_N(z^2),\qquad \deg g_N=M.
\end{equation}
The degree is exactly $M$: otherwise the nonzero raising coefficient
$(N-k)(N-k-1)$ at the highest even degree $k<2M$ would contradict the
eigenvalue equation. The equation for $g_N$ is
\begin{equation}\label{eq:g-ode}
 \begin{aligned}
 4t(t^2-1)g_N''
 &+\bigl((6-4N)t^2-4(N+1)t-2\bigr)g_N'\\
 &+\bigl(N(N-1)t+N(N+1)-\tau_N\bigr)g_N=0.
 \end{aligned}
\end{equation}
The other parity sector is obtained by substituting $F(z)=z g(z^2)$
in \eqref{eq:H}; it is not the selected Perron sector used here.

\section{The compact-interval Heun polynomial}
Reverse the polynomial by
\begin{equation}\label{eq:y-def}
 y_N(x)=\frac{x^Mg_N(-1/x)}{(-1)^M[t^M]g_N(t)}.
\end{equation}
Then $y_N(0)=1$. Define
\begin{equation}\label{eq:q}
 \gamma=d+\frac12,\qquad
 q_N^*=\frac{\tau_N+(N+1)(N-2d)}4.
\end{equation}
Substitution in \eqref{eq:g-ode} gives
\begin{equation}\label{eq:heun}
 \boxed{y_N''+
 \left(\frac{\gamma}{x}+\frac1{x-1}-\frac N{x+1}\right)y_N'
 +\frac{M(M-\frac12)x-q_N^*}{x(x-1)(x+1)}y_N=0.}
\end{equation}
In the convention
\[
 y''+\left(\frac\gamma x+\frac\delta{x-1}+\frac\epsilon{x-a}\right)y'
 +\frac{\alpha\beta x-q}{x(x-1)(x-a)}y=0,
 \quad \alpha+\beta+1=\gamma+\delta+\epsilon,
\]
our parameters are
\[
 a=-1,\quad \alpha=-M,\quad\beta=\frac12-M,\quad
 \gamma=d+\frac12,\quad\delta=1,\quad\epsilon=-N,\quad q=q_N^*.
\]
Thus the analytic branch normalized at zero can be written
\[
 y_N(x)=\HeunG\left(-1,q_N^*;-M,\frac12-M,d+\frac12,1;x\right).
\]
The selected accessory parameter enforces polynomial termination and
selects the highest polynomial mode; termination alone has $M+1$
possible accessory parameters.

\subsection{Coefficient recurrence}
Writing $y_N(x)=\sum_{n=0}^Mc_nx^n$, with $c_{-1}=0,c_0=1$, gives
\begin{equation}\label{eq:rec}
 c_{n+1}=
 \frac{((N+1)n-q)c_n+(n-M-1)(n-M-\frac12)c_{n-1}}
 {(n+1)(n+\gamma)}.
\end{equation}
In particular, $c_1=-q/\gamma$. Termination is $c_{M+1}(q)=0$.
This accessory polynomial is distinct from the spectral-return
polynomials $p_N$ and $\widetilde q_N$ used elsewhere in the project.

\section{Positive-weight Sturm--Liouville form}
On $0<x<1$ define
\[
 w_N(x)=x^{d-1/2}(1+x)^{-N-1},\qquad
 a_N(x)=x^{d+1/2}(1-x)(1+x)^{-N}.
\]
Equation \eqref{eq:heun} is equivalent to
\begin{equation}\label{eq:SL}
 -\frac1{w_N}(a_Ny_N')'
 +M\left(M-\frac12\right)xy_N=q_N^*y_N.
\end{equation}
The weight is positive and integrable. Natural zero-flux conditions
$\lim a_Ny'=0$ at both endpoints select the regular polynomial branch.
The differential expression preserves $\Poly_M$ and is symmetric there
in $L^2((0,1),w_Ndx)$.

The selected $q_N^*$ is the largest eigenvalue in this finite invariant
polynomial space, not the largest eigenvalue of the full unbounded
Sturm--Liouville problem. The polynomial eigenfunctions occupy its
first $M+1$ natural-boundary modes: each has at most $M$ zeros, while
Sturm oscillation assigns distinct zero counts to distinct eigenvalues.
Consequently the top polynomial has $M$ simple zeros in $(0,1)$:
\begin{equation}\label{eq:zeros}
 0<x_{1,N}<\cdots<x_{M,N}<1.
\end{equation}
For $M=0$, all root products below are empty products.

\section{Returning to the original symmetric matrix}
Write $\lambda_N=\lambda_{\max}(S_N)$; the notation $L_N$ remains
reserved for the lowering operator.
The exact return identities are
\begin{equation}\label{eq:return}
 \boxed{\lambda_N
 =2^{N/2}\left|\frac{g_N(1)}{g_N(-1)}\right|
 =2^{N/2}\left|\frac{y_N(-1)}{y_N(1)}\right|.}
\end{equation}
For an explicit check, $S_Nf_N=c_N\sigma f_N$ and $\Phi_N^2=2^N\Id$
imply, after applying $\Phi_N$ and taking the constant coefficient,
\[
 2^Nf_N(0)=c_Nf_N(-1).
\]
On the other hand, \eqref{eq:C2} gives
$g_N(-1)=F_N(i)=(2i)^Nf_N(-1)$ and
$g_N(1)=F_N(1)=[2(1+i)]^Nf_N(0)$.
Taking absolute values proves the first equality in \eqref{eq:return};
\eqref{eq:y-def} proves the second.
Equivalently, using \eqref{eq:zeros},
\begin{equation}
 \lambda_N=2^{N/2}\prod_{j=1}^M\frac{1+x_{j,N}}{1-x_{j,N}}.
\end{equation}
The coefficient of the next-to-highest power in \eqref{eq:heun} gives
\begin{equation}
 \tau_N=N(N+1)-2\sum_{j=1}^Mx_{j,N}.
\end{equation}
These are exact finite identities. They do not, by themselves, justify
weak zero-measure limits, endpoint logarithmic statistics, or a sharp
large-$N$ prefactor.

\section{A second Heun chart with third singularity at two}
For reference, set
\begin{equation}\label{eq:U}
 s=\frac{2x}{1+x},\qquad
 U_N(s)=\left(1-\frac s2\right)^M
 y_N\!\left(\frac{s}{2-s}\right).
\end{equation}
This is a degree-$M$ polynomial with $U_N(0)=1$ and Heun parameters
\begin{equation}\label{eq:U-par}
 \begin{gathered}
 a=2,\quad Q_N=-q_N^*-M\gamma,\quad
 \alpha=-M,\quad\beta=M+d+1,\\
 \gamma=d+\frac12,\quad\delta=1,\quad\epsilon=\frac12.
 \end{gathered}
\end{equation}
Thus $U_N''+(\gamma/s+1/(s-1)+1/(2(s-2)))U_N'
 +(\alpha\beta s-Q_N)U_N/[s(s-1)(s-2)]=0$.
Moreover,
\[
 U_N(1)=2^{-M}y_N(1),\qquad
 [s^M]U_N=(-1)^M2^{-M}y_N(-1),
\]
so
\begin{equation}
 \lambda_N=2^{N/2}\left|\frac{[s^M]U_N}{U_N(1)}\right|.
\end{equation}
The third finite singularity is now outside the unit disk. The
exponent difference at $s=1$ is nevertheless zero, and the difference
at infinity is $N+1$. Generic nonresonant connection formulas therefore
require a controlled limiting procedure. An expansion in $1/a=1/2$
is not automatically uniform as the other parameters grow with $N$.

\section{Relation to the current asymptotic approach}
The differential route is an exact alternative description, not a
necessary part of the present leading-asymptotic strategy. If $v_N$ is
the Jacobi Perron vector (with any normalization), then
\eqref{eq:normbridge} yields directly
\begin{equation}\label{eq:positive-ratio}
 \boxed{\lambda_N^2=2^N
 \frac{\sum_{j=0}^N\binom Nj(v_N)_j^2}
 {\sum_{j=0}^N\binom Nj^{-1}(v_N)_j^2}.}
\end{equation}
The expected central profile, with $\xi_j=(j-N/2)/\sqrt N$, is
$(v_N)_j\propto e^{-\sqrt8\,\xi_j^2/2}$; a convenient discrete
candidate is $(v_N)_j\propto\binom Nj^{1/\sqrt2}$.
Local Gaussian convergence plus suitable inverse-binomial weighted tail
control would give
\[
 \lambda_N\sim\sqrt{\frac{\sqrt8-2}{\pi}}\,
 \frac{(\sqrt8)^N}{\sqrt N}.
\]
That weighted control is a separate proof obligation; no such
asymptotic theorem is asserted here.

\section*{Notation map}
\begin{center}
\begin{tabular}{ll}
$f_N(z)$ & Original $T_N$ polynomial eigenfunction, $f_N(0)=1$.\\
$F_N(z)$ & Cubic conjugate $C_N^2f_N$, even for the top mode.\\
$g_N(t)$ & Half-degree polynomial, $F_N(z)=g_N(z^2)$.\\
$y_N(x)$ & Reversed Heun polynomial, $y_N(0)=1$.\\
$U_N(s)$ & Alternate Heun polynomial with $a=2$.\\
$\tau_N$ & Largest eigenvalue of the positive Jacobi matrix.\\
$q_N^*$ & Selected Heun accessory parameter in the $x$ chart.\\
$Q_N$ & Accessory parameter in the $s$ chart.\\
$\lambda_N$ & Largest eigenvalue of the symmetric matrix $S_N$.\\
$\sigma,\cI$ & Parity and phase matrices; $\cI^2=\sigma$.
\end{tabular}
\end{center}
\end{document}
